\par\bigskip
\begin{center}
\rule{.88\textwidth}{.3pt}\par
Göttingische gelehrte Anzeigen.\quad 25 September 1826.\par
\rule{.88\textwidth}{.3pt}
\end{center}
\pdfbookmark[1]{25 September 1826}{anzeige-1826-09-25}

On 16 September Councillor \textsc{Gauss} presented to the Royal Society a lecture entitled:
\begin{center}
\textit{Supplementum Theoriae combinationis observationum erroribus minimis obnoxiae.}
\end{center}
\noindent In all earlier works on applying the calculus of probability to the most suitable use of observations, and particularly in the treatment of this subject in the fifth volume of the \textit{Commentationes recentiores}, the form of the main problem rests on a definite assumption that is indeed suited to most cases encountered in practice. This assumption is that the ob\SourceLeaf{110}%
served quantities depend in a known way on certain unknown quantities (elements), that is, that they are known functions of these elements. For the problem to belong to this class at all, the number of these elements must be smaller than the number of observed quantities, so that the latter are themselves dependent on one another.

Nevertheless, the assumed condition is not directly met in many cases: the observed quantities are not yet given as known functions of certain unknown elements, and it is not immediately apparent how to put them in that form. Instead, their mutual dependence—which must, of course, be given in some way—is expressed by certain equations of condition, which, by the nature of the matter, the true values of those quantities must satisfy exactly. On closer examination one soon sees that this case is not essentially different from the other, but differs only in form, and in theory can indeed readily be reduced to it. In application, however, this is often an unnatural detour that leads to calculations far more troublesome than a solution of its own, specially suited to the original form of the problem. The present treatise is devoted to such a solution. The solution of the problem, which the treatise gives as a self-contained whole independent of the earlier treatise, is itself so flexible that in some cases it may even be advantageous to use it where the assumption underlying the older method already holds.

The main problem now takes the following form. If another quantity depends in some way on the quantities \(v,\,v',\,v''\), and so on, among which a relation holds that is given by one or more equations of condition—for example, if it can be expressed by the function \(u\)—then it can equally be determined in infinitely many other ways from those quantities, or expressed by infinitely many other functions in place of \(u\). Naturally, all these functions give the same results when the true values of \(v,\,v',\,v''\), and so on, which satisfy all the equations of condition, are substituted. But if one has only approximate values of \(v,\,v',\,v''\), and so on, as observations of limited accuracy can always provide, then \SourceLeaf{111}%
the quantities derived from them cannot claim absolute correctness either. Generally speaking, the different functions used for \(u\) will give unequal results, but, more importantly, results of unequal reliability. The problem is to select, from the infinite variety of functions by which the unknown quantity can be expressed, the one whose result is liable to the smallest possible unreliability.

The treatise in fact gives \textit{two} solutions to this problem. The first reaches the goal by the shortest route when only one unknown quantity, depending on the observations in a prescribed way, is actually to be derived. But closer consideration of this solution also leads to the remarkable theorem that the very same value for the unknown quantity as follows from the most suitable combination of the observations is obtained if one first makes certain changes in the observations, calculated according to specified rules, and then substitutes them into any arbitrary function expressing the unknown quantity. Besides satisfying all the equations of condition, these changes have the property that, among all conceivable systems that do so, the sum of their squares (assuming the observations to be equally reliable) is as small as possible. Thus a new foundation is obtained for the method of least squares, and this \textit{adjustment} of the observations, wholly independent of the function \(u\), furnishes a second solution that has a great advantage over the first when more than one unknown quantity is to be derived from the observations in the most suitable way: the observations are thereby prepared in advance for \textit{any} use to be made of them. This second solution still required special instructions for determining the degree of accuracy achieved in each individual application. The treatise contains complete instructions for all this, made as simple as possible; naturally, they cannot be excerpted here. Nor can we go into the details here of the further investigations carried out after the development of the main problems, although they are closely connected with the subject. We mention only one remarkable theorem: \SourceLeaf{112}%
the rules for the complete adjustment of the observations always give the same result whether applied to the original observations themselves or to observations that have already been \textit{incompletely} adjusted for the time being, provided that this concept is understood in the more precise sense defined in the treatise. As a special case, this includes the situation in which the observations have already undergone an adjustment carried out according to the rules but taking account of only part of the equations of condition.

The final part of the treatise consists of a couple of carefully worked examples of the method’s application, drawn in part from the geodetic measurements of General von \textsc{Krayenhoff} and in part from the triangulation carried out by the author himself in the Kingdom of Hanover. They can serve both to explain this theory’s application further and to shed more light generally on circumstances relating to measurements of this kind.

Trigonometric measurements belong especially to the field in which the calculus of probability applies, and particularly to the form of its application developed in the present treatise. Here it is the rule to make more observations than are strictly necessary, so that the measurements provide many mutual checks. Only by using the strict principles of the calculus of probability can one fully obtain the benefit available from this circumstance and give the results the greatest accuracy of which they are capable. These principles also provide the means to determine the accuracy of the measurements themselves and the admissibility of the results based on them. Finally, they make it possible to select the most suitable arrangement of a system of triangles from several among which one may have a choice. All this can be done according to firm and reliable rules, excluding all arbitrariness. Yet measurements can be judged with confidence and used to the fullest extent only when they are presented in unaltered authenticity and completeness. It would therefore be very desirable for all larger measurements of this kind that claim exceptional accuracy always to be published with all necessary detail. Far too often the opposite occurs, with only the final results for the individual mea\SourceLeaf{113}%
sured angles being communicated. If such final results are formed according to sound principles, by allowing all individual series of observations to contribute except those that certainly contain a wholly inadmissible error, the disadvantage is indeed far less than when one keeps only those series that best fit the nearby checks afforded by the sums of the angles in each triangle and by the sums of the horizontal angles around each point. When this wholly reprehensible procedure is used, whether from ignorance of the true principles of a sound theory or from a secret wish to give the measurements the appearance of greater accuracy, the standard for a fair assessment of the observations and the results to be derived from them is lost. The usual check by angle sums in individual triangles, and at points where the measured angles cover the whole horizon, then seems to prove a measurement accuracy that may in fact be very far from being attained; and if other checks are available, through side ratios in closed polygons or through diagonal directions, these will reveal with certainty the existence of much larger errors. Conversely, when the condition just mentioned holds and the observations are adjusted with respect to the checks without the reliable rules of the calculus of probability, the process must remain a groping in the dark and leads to larger, often much larger, corrections than necessary. This can readily give rise to an \textit{unduly} unfavorable judgment of the measurements. These remarks show the importance both of publishing measurements in sufficient detail and of mathematically combining geodetic measurements on strict principles. But they evidently apply, to a greater or lesser extent, to observations of every kind—astronomical, physical, and so on—that concern quantitative matters, insofar as the variety of circumstances involved affords means of mutual checks.
\par\medskip\centerline{\rule{.3\linewidth}{.3pt}}
